\documentclass[11pt]{article}
\usepackage{amsmath,amssymb,hyperref}
\title{From resilient corridor formation to timely report delivery}
\author{}\date{}
\begin{document}\maketitle
Consider ordered moving radios on a line between two fixed, loss-exempt
anchors. A stopped body remains a fixed physical obstacle (not a radio).
The two sides know its position, and each side knows its own positions
and speed limits. All paths must maintain minimum clearance $g>0$.
A schedule of duration $T$ constrains each endpoint to its reachable
interval. The distance-only symmetric graph uses range $R$, and the
terminal guarantee allows up to $k$ additional interior radios to go
silent without moving their bodies. Physical clearance is continuous;
communication here is required after motion freezes, not throughout
approach. A single retained report must go from one anchor to the other.

\paragraph{Terminal connectivity.}
In the unfailed sequence $x$, test $x_j-x_i\le R$ whenever
$1\le j-i\le k+1$. If a test fails, deleting the intervening radios
disconnects the two sides of that span; they are then consecutive,
and no path can cross the gap. Conversely, consecutive survivors have
at most $k$ deleted radios between them and their span was tested.
On each side the physical bounds, speed intervals, and robust span
constraints are difference bounds. Each local feasible set is closed
under coordinatewise minimum and maximum. For an upper inequality
$x_j\le x_i+c$, both endpoint choices are preserved by each lattice
operation. Since each side is closed and bounded, combining extremizers
yields its componentwise greatest and least feasible vectors. Choose the
greatest vector on the left and least on the right. Those choices improve
every cross-body reach bound, and the known body gives cross clearance.
The affine path from the current state to the chosen endpoints preserves
the convex clearance conditions and rate bounds. Feasible terminal
connectivity is therefore exact when the two locally chosen vectors
meet the cross tests. A negative result is a statement about this
synchronized, known-obstacle model, not a recovery rule after another
physical obstruction appears. The comparison is a direct full-system
difference-constraint solver under the same information.

This computation does not itself verify the terminal radio edges. In
a one-bit query-only model there is a sharp example for individual cross
tests. Use nominal unit spacing, range $q+3\epsilon/2$, $q=k+1$, and
perturb one nonanchor cross span of index length $q$ by moving its left
endpoint left by $\epsilon$ and its right right by $\epsilon$. For small
$\epsilon$ this respects the same known obstacle clearance. The target
span now fails while all others of length at most $q$ pass. The all-pass
history is indistinguishable without that particular query. The lower
bound does not apply to exact position data or analog measurements.

\paragraph{A service deadline is stronger than connectivity.}
With static frozen failures and unit-hop travel, the farthest surviving
forward neighbor minimizes hops. For any competing path, the greedy
position at each hop is at least as far along, by monotonicity of the
rightmost reachable coordinate. To maximize greedy hop count under
deletions, store the current index $i$, the former tested frontier $e$
so $(i,e]$ is already silent, and the remaining deletion budget. Make
$j>e$ the next greedy neighbor by silencing the suffix above it in the
new reachable prefix. The adversary may defer all other deletions until
they are first needed, eliminating hidden state beyond the frontier.
Killing all reachable fresh neighbors disconnects; the end anchor cannot
be killed. This gives the recurrence in the replay and matches explicit
deletion plus breadth-first search. Static unbounded-hop connectivity
alone cannot certify a report deadline.

For a two-hop deadline, if the anchors are not direct neighbors, the
intermediate radio must lie in
$[x_{\mathrm{right}}-R,x_{\mathrm{left}}+R]$. It follows exactly that at
least $k+1$ interior radios must occupy this interval. They form a
consecutive block, so enumerate candidate block starts and impose
simple position bounds on those radios; use the split solver subject
to those additional bounds. Block choice is a shared decision in this
construction, not a claim that both sides can choose any adequate block
in a blackout. This gives a continuous-deadline test, because the motion
and spacing constraints remain convex in each retained block. Same-cut
verified anchor ranges on ordered identified radios suffice to count
the intersection of the left prefix and right suffix. Reports from
different cuts alone do not: a single relay can move from 4 to 7 between
observations for anchors 0 and 11 and range 6. It was heard by the left
anchor and later by the right, but is not a two-hop relay at the right
anchor's cut. A common phase or a true dwell guarantee is needed.

\paragraph{Strict geometric example.}
Let the starting vector be $(0,2,4,6,7,8,10,11)$, the body at 3, speed
limits $(0,0,1,0,1,1,1,0)$, $g=1$, $R=6$, and $k=1$. Robust connectivity
already holds at time zero. At $T=1$, cut-directed extrema give
$(0,2,4,6,7,8,9,11)$, still requiring three hops in the worst failure.
The two-hop interval is $[5,6]$. To place two radios there one must
move the initially-4 radio to 5, which takes at least one unit; the
next radio is fixed at 6 and the remaining right radios cannot move
past it or into its clearance margin. The placement
$(0,2,5,6,7,8,9,11)$ realizes the exact deadline and retains robust
two-hop service. Stronger service changes the movement and evidence
interfaces even with identical range, safety and speed constraints.

\paragraph{Actual custody and failure timing.}
The static route calculation is not a delivery theorem if loss can
occur after the first hop. In a separate two-epoch model, the sender can
actually store at most $a$ durable copies in distinct intermediates of the
two-hop interval, and up to $k$ can go silent after storage and before
forwarding. There are at most $c$ collision-free forwarding slots in
the second epoch, one assigned relay per slot. Anchors survive and there
is a single report. With no timely identification of failed relays, the
preassigned eligible relays that actually hold a copy must number at
least $k+1$. If $n$ radios are in the interval, this is possible exactly
when $\min(n,a,c)>k$. With timely truthful identification before the
slot assignment, this instead requires $\min(n,a)>k$ and $c\ge1$.
Necessity follows by deleting all eligible candidates when there are
no more than $k$; sufficiency delivers on a slot from a surviving actual
copy. This counts real custody and scheduling capacity. Neither a
certificate for the initial graph nor receipt of an instruction creates
a copy, time slot, or liveness under a later loss. These are supplied
source primitives, not guarantees about a specific wireless medium.

\paragraph{Comparison.}
\href{https://publications.ri.cmu.edu/minimum-k-connectivity-maintenance-for-robust-multi-robot-systems}{Luo and Sycara}
study continuous robust connectivity maintenance.
\href{https://arxiv.org/abs/2404.03834}{Ishat-E-Rabban et al.}
study general movement-minimizing connectivity restoration with a QCP
and an approximation-oriented search. This reference specializes to
an ordered corridor with a fixed obstacle and distinguishes exact
terminal recovery from a report service deadline. The full difference
solver, explicit deletion/BFS and actual slot checks are matched
comparators; no speed or priority claim against those research methods
follows. Finite replay covers 106,620 deletion cases, 10,000 rational
placement comparisons, 715 necessary cross-probe witnesses, 21,570
hop checks, 2,000 service placement checks, 21,555 anchor-evidence checks,
and 1,250 custody/slot checks. Five Lean lemmas check supporting
arithmetic, lattice, and cardinality statements.
\end{document}
